% =====================================================================
%  Interface Friction as a Constraint on Market Participation
%  Taylor & Francis / Interact class (JOCEC, hoce20).
%
%  Rendered against: paper/template/journal/main.tex  (interact v1.05)
%  Source of truth:  paper/paper-3-jocec-barrier-taxonomy.md
%  Bibliography:     paper/paper-3-jocec-references.bib
%  Plan:             paper/PLAN-3-paper-split.md
%
%  Build:  pdflatex paper-3-jocec-barrier-taxonomy.tex
%          bibtex   paper-3-jocec-barrier-taxonomy
%          pdflatex paper-3-jocec-barrier-taxonomy.tex   (x2)
%
%  NOTE: this file has NOT been compiled -- no LaTeX toolchain is
%  installed on this machine. See the handoff note in the repository.
% =====================================================================
\documentclass[]{interact}

\usepackage{epstopdf}
\usepackage[caption=false]{subfig}
\usepackage{booktabs}
\usepackage{longtable}
\usepackage{amsmath}
\usepackage{amssymb}
\usepackage[longnamesfirst,sort]{natbib}
\bibpunct[, ]{(}{)}{;}{a}{,}{,}
\renewcommand\bibfont{\fontsize{10}{12}\selectfont}

\theoremstyle{plain}
\newtheorem{proposition}[section]{Proposition}

\begin{document}

\articletype{ARTICLE}

\title{Interface Friction as a Constraint on Market Participation: A Typology and Remediation Matrix for MSME Service Marketplaces}

\author{
\name{Bagaskoro Saputro\textsuperscript{a}\thanks{Corresponding author. Email: bagaskoro.saputro@binus.ac.id} and Yashella Tirana\textsuperscript{b} and Freysia Chandra Saliman\textsuperscript{c} and Laurentius Ryan Dofani Pradita\textsuperscript{c} and Gregory Adrianus Sugiono\textsuperscript{c} and Gregorius Andrew Rahardja\textsuperscript{c} and Steven Karunia, Liem\textsuperscript{c} and Joshua Joyson Kustiadi\textsuperscript{c}}
\affil{\textsuperscript{a}BINUS School of Computer Science, Computer Science, Jakarta, Indonesia; \textsuperscript{b}BINUS Business School, Digital Business, Jakarta, Indonesia; \textsuperscript{c}BINUS School of Computer Science, Computer Science, Semarang, Indonesia;}
}

\maketitle

\begin{abstract}
Digital platforms are widely promoted as the route by which micro, small and medium enterprises
(MSMEs) reach customers beyond their physical catchment. The premise is contested in this paper. We
report a formative evaluation of an Indonesian-language marketplace prototype built for home-maintenance
firms, and we find that \textbf{interface friction is not a cosmetic layer over digitalisation but a
constraint on market participation itself.} Ten participants completed eight standardised tasks across
two moderated rounds (80 task attempts), cross-validated against an independent code-based heuristic
audit of the same prototype.

Six barrier themes emerge, which we organise into a \textbf{typology of five friction classes} --- service-concept
ambiguity, infrastructure asymmetry, representational inconsistency, transaction-state fragility, and
support-channel occlusion. Each class is grounded in a construct from organisational computing or
information systems research, and each is stated as a proposition with boundary conditions (P1--P5).
Two findings warrant attention beyond this prototype. First, a matched pair of discovery tasks
isolates a \emph{taxonomy} defect from a \emph{layout} defect --- the same task, differing only in service
category, varies by 0.60 SEQ points and triples in variance, indicating a low-cost data remedy rather
than an interface redesign. Second, and more consequential for practice, the address form reached the
ceiling of the ease scale while generating the most-requested features in the entire study:
\textbf{ease of completion and fitness for purpose are independent properties, and the common post-task
instrument measures only the first.}

We contribute a friction typology, a remedy-cost classification (data change versus interface change
versus infrastructure provision), and an evidence-linked remediation matrix in which every recommended
action is traceable to a frequency count, a dispersion measure, and an independent code finding. We
report that the narrative source report contradicts its own prioritisation of support-channel
remediation, and we resolve the contradiction in favour of the recorded evidence rather than the
report's own conclusion. Findings are bounded by five disclosed provenance defects.
\end{abstract}

\begin{keywords}
MSME digitalization; e-commerce platform; interface friction; service marketplace; organisational
computing; technology--organisation fit; formative evaluation; Indonesia
\end{keywords}

\section{Introduction}

\subsection{The digitalisation premise for Indonesian MSME service firms}

Household maintenance services in Indonesia --- air-conditioner cleaning, plumbing leak repair, pump
servicing, septic tank suction --- are delivered predominantly by micro and small enterprises. Their
customer acquisition still runs through \emph{spanduk} (street banners), printed leaflets, and vehicle
stickers. This channel has three structural properties, each of which is invisible to the firm as well
as to the customer:

\begin{enumerate}
\item \textbf{A fixed physical area}, which caps how much information can be presented at once.
\item \textbf{No record of prior work}, which is replaced by informal reputation.
\item \textbf{No displayed price}, because the price is negotiated privately and is therefore
      unverifiable in advance.
\end{enumerate}

The third property is the commercially decisive one. A privately negotiated price cannot be compared,
cannot be budgeted, and cannot be checked for overcharging, because there is no second number to check
it against. The firm has no channel through which to signal reliability, and the customer has no basis
on which to select.

Digital platforms are proposed as the remedy. \textbf{We take the premise as given and test what it
costs.} Our question is not whether platforms help MSME firms reach customers --- a large literature
assumes they do --- but what a platform interface does to a service firm's \emph{capacity to transact
credibly at scale}, and what happens when it fails to.

\subsection{Why interface friction is an organisational problem, not a usability problem}

The prevailing framing of this material treats interface defects as usability problems: local, fixable,
and exhausted by a list of corrections. We argue this framing systematically understates them in MSME
service marketplaces, for a reason specific to that context.

A marketplace interface is the \textbf{firm's principal asset for demonstrating competence to a
stranger}. In a transaction where the firm sends an unknown person into a customer's private home, the
interface is the only place where the firm can establish that it is organised, legible, and
accountable --- before any interaction has occurred. Interface friction is therefore not friction in
the user's journey; it is friction in the \textbf{firm's capacity to be trusted by someone who has
never heard of it}.

This has three organisational consequences:

\begin{itemize}
\item \textbf{Trust is pre-emptive and non-recoverable.} A customer who encounters two contradictory
      prices, or a help channel that reads as a booking action, acquires a reason to distrust every
      subsequent number. There is no recovery path within the interface.
\item \textbf{Remediation cost is not proportional to user effort.} Several of the barriers we identify
      are remediable by a \emph{data change} --- merging two catalogue entries --- rather than by a
      redesign. Framing these as usability defects misallocates scarce developer time toward expensive
      fixes.
\item \textbf{The firm absorbs the cost, not the customer.} In this transaction the customer cannot
      substitute. Choosing another provider means restarting the search from zero, which is why the
      interface's failure to communicate support converts a design flaw into an abandoned or
      unsupported booking.
\end{itemize}

We therefore treat the interface as an organisational instrument and the barriers as a \textbf{typology
of friction in market participation}.

\subsection{Research questions}

\begin{itemize}
\item \textbf{RQ1.} What classes of interface friction arise when an MSME service firm moves its
      transaction onto a marketplace platform?
\item \textbf{RQ2.} Which barriers are remediable by a data change, which by an interface change, and
      which require infrastructure the platform must provide?
\item \textbf{RQ3.} Where do participant observation and independent code inspection converge, and where
      does each uniquely contribute?
\item \textbf{RQ4.} How should a formative evaluation prioritise remediation when the source records
      disagree with one another?
\end{itemize}

\subsection{Contributions}

\begin{enumerate}
\item \textbf{A five-class typology of interface friction in MSME service marketplaces} (service-concept
      ambiguity, infrastructure asymmetry, representational inconsistency, transaction-state fragility,
      support-channel occlusion), each grounded in an established IS or organisational-computing
      construct and stated as a proposition with boundary conditions.
\item \textbf{A remedy-cost classification.} Data change versus interface change versus infrastructure
      provision --- a distinction that changes remediation budgeting and is invisible in standard
      usability reporting.
\item \textbf{A matched-pair diagnostic} that attributes a difficulty to a taxonomy defect rather than a
      layout defect, converting a vague complaint into a specific and cheap remedy.
\item \textbf{An evidence-linked remediation matrix} in which every recommended action traces to a
      frequency count, a dispersion measure, and an independent code finding, so that priority is
      auditable rather than asserted.
\item \textbf{A documented resolution of a source contradiction}, reported rather than suppressed: the
      narrative source's own recommendation matrix contradicts its own session data on support-channel
      remediation.
\end{enumerate}

\subsection{Scope}

This is a formative evaluation of a single prototype build. No variable measuring conversion, revenue,
bounce rate, cognitive load, or trust change was instrumented, and \textbf{we make no claim about any
of them}. In particular, we do not claim that the interface caused any change in firm revenue or
customer acquisition; we have no measurement of either, and a single-prototype formative study cannot
support such a claim. Our claims are about the \emph{characterisation and prioritisation of friction},
which is evidenced by what participants did and by what the code does.

\section{Related work}

\subsection{Digitalisation of MSME firms and technology--organisation fit}

The premise that digital platforms extend MSME market reach is well established. Tianti et al.\
\citep{tianti2024} demonstrate that a small catering firm's website increased promotional reach;
Sunarso et al.~\citep{sunarso2023} provide evidence that visual content raises attraction and
conversion in MSME digital marketing. We adopt the premise and decline to import effect sizes measured
on different stimuli and markets.

Technology--organisation fit research \citep{nadler1989,tbdTechnologyOrganisationFit} supplies the
analytical frame this paper adopts. Fit is usually described as congruence between a technology's
capabilities and an organisation's requirements. We invert the question. In an MSME service marketplace
the fit problem is not primarily whether the platform's \emph{capabilities} match the firm's
\emph{needs}, but whether the interface's \emph{representation} matches what the firm can stand behind
--- its service taxonomy, its prices, its operating procedures. A platform can offer every capability a
firm needs and still misrepresent the firm's structure, and our data show a case in which it does.

\subsection{Platform affordances and infrastructure asymmetry}

Platform research distinguishes capabilities a platform offers from what a user or firm can actually
exploit. The affordances literature \citep{norman2013} identifies how an interface's properties guide
action; technology capability research \citep{tbdPlatformCapability} distinguishes what a system enables
from what it delivers. Our second friction class concerns what the platform does \emph{not} offer.
Participants repeatedly requested geocoding, address autocomplete, saved addresses, order history, and
re-order. Collectively these describe a checkout that asks the firm to have its customer re-supply known
information on every booking.

We treat this as \textbf{infrastructure asymmetry}: work that the platform could perform once and reuse
is delegated to the user on every transaction. For a micro-firm whose entire advantage is repeat
custom, this inverts the intended economics of the platform.

\subsection{Information consistency and the representation of organisational capability}

Organisational computing research long distinguished between what an information system
\emph{stores} and what it \emph{represents} as the organisation's activity. Our third friction class
concerns consistency between representations: two date pickers offering disagreeing options, a
home-page price differing from the catalogue, a tax label whose stated rate does not determine its
displayed amount \citep{braniwati2026price}.

Consistency here is not aesthetic. A redundant surface is harmless when the copies agree and actively
harmful when they do not, because the user can no longer determine which is authoritative. In a
pricing context the user resolves the ambiguity by trusting neither. A companion paper
\citep{braniwati2026price} develops this argument at length, grounds it in metamorphic testing, and
reports a reproduction attempt that refuted its own headline finding --- we draw the consistency
principle from it and cite it rather than restating it.

\subsection{Formative evaluation instruments and their limits at task level}

Single Ease Question (SEQ) \citep{sauro2016} is a widely adopted post-task instrument: one 7-point item
administered immediately after each attempt. It yields breadth --- one comparable number per attempt
--- at the cost of discarding within-task structure.

Our companion paper \citep{braniwati2026dispersion} reports the consequence sharply. In this same
study, \textbf{four of eight tasks saturated at $SD \leq 0.42$}, meaning they could no longer
distinguish a good design from a very good one. Had we ranked remediation priority by raw frequency
across all barriers, we would have weighted saturated tasks --- which generated no discriminative
evidence --- equally with the one task that spanned the full scale. We therefore use dispersion as a
\textbf{gate} on which barriers enter the typology, and frequency counts only among tasks that produced
variance. This is a methodological choice with consequences for the backlog, and we state it explicitly
because the alternative is defensible and would have produced a different priority ordering.

\subsection{Heuristic inspection as a complementary instrument}

Heuristic inspection \citep{nielsen1994} and its Indonesian applications
\citep{perdanakusuma2022,febriyanthi2021} inspect code against a model of intended use. Our prototype
received an independent code-based audit producing 0 Critical, 2 Major and 4 Minor findings. The
division of labour between the two methods is quantified in Section~\ref{sec:division}.

\subsection{Design rationale as background, not as measured effect}

Our prototype drew on Norman's three levels of design \citep{norman2013}, Hick's Law for navigation
depth \citep{hick1952}, Miller's Law for grouping \citep{miller1956}, and Von Restorff isolation for
action prominence \citep{vonrestorff1933}. These function as \emph{rationale}. None was manipulated as
an independent variable, and \textbf{we make no claim that any principle caused any observed
outcome.} We cite them to describe the design intent against which participants' behaviour is
interpreted, which is the only claim the study design can support.

\section{Method}
\label{sec:method}

\subsection{Prototype under test}

A high-fidelity React prototype, localised into Indonesian, for on-demand household maintenance,
implementing service discovery with search and category filtering, service detail pages, a four-stage
checkout (summary $\rightarrow$ schedule $\rightarrow$ address $\rightarrow$ payment), a
confirmation/invoice view, technician tracking, and an FAQ/contact area.

Two properties were deliberately left unchanged because they bear directly on our research questions:
prices render in \textbf{USD} while the interface is Indonesian-language, and service descriptions
include \textbf{technician-internal procedure content} intended for engineers rather than customers.

\subsection{Participants and disclosure}

Ten participants were recruited and tested in two rounds of five, identified only by anonymous codes
(P01--P05, P06--P10). Sessions were conducted in Semarang, Central Java, between \textbf{1 and 30
October 2026}.

\textbf{Disclosure on demographics.} Participants were Gen Z, aged \textbf{19--21}, verified against
their national ID cards by the research team at recruitment and confirmed for all ten. No date of
birth was transcribed into the research data, as a privacy measure; the age range is recorded here as an
aggregate. Gender and occupation were not collected. We describe the sample as \emph{Gen Z users
recruited purposively} and make \textbf{no representativeness claim} beyond that description.

\subsection{Tasks and measures}

Eight tasks spanning the booking flow. T1 and T2 are a \emph{matched pair} differing only in service
category. T6 and T7 are \textbf{defect probes} with explicit failure criteria. Measures: task outcome
(completed / completed with assistance / not completed against a written criterion), SEQ on a 7-point
scale administered immediately after each attempt, and observer commentary.

Sessions were moderated with concurrent think-aloud. Screen and audio recording was conducted with the
consent of all ten participants and is stored under institutional control, scheduled for deletion on
1 January 2027; \textbf{the recordings were not used in this analysis}, and all qualitative material is
observer paraphrase, not verbatim transcript. We use quotation marks only for source documents, never
for participant speech.

\subsection{The completion-rate dispute, and how we resolved it}

Completion was coded during sessions and recorded as 39/40 (round 1) and 40/40 (round 2). Both source
recap sheets then appended an auditor note that specific attempts should have been coded failures:
round 1, two T6 attempts recorded as successes although the criterion was \emph{find the validation
message} and the notes state the participants did not find it; round 2, one T6 attempt recorded as a
success although the note reads ``user did not find any error'', and another reaching confirmation with
a street address consisting only of the characters \texttt{(} and \texttt{)}.

This is an interpretive disagreement about whether failure on a \emph{sub-criterion} constitutes task
failure, not an arithmetic error. \textbf{We report the auditor-corrected figures as the headline and
present the as-recorded figures as sensitivity analysis.}

\begin{table}[htbp]
\caption{Task success, both codings.}
\begin{tabular}{@{}lccc@{}}
\toprule
 & Round 1 & Round 2 & Combined \\
\midrule
As recorded & 39/40 (97.5\%) & 40/40 (100\%) & 79/80 (\textbf{98.75\%}) \\
Auditor-corrected (headline) & 37/40 (92.5\%) & 38/40 (95.0\%) & 75/80 (\textbf{93.75\%}) \\
\bottomrule
\end{tabular}
\end{table}

All five failures occurred at T6 or T8.

\subsection{How barriers entered the typology}

Two filters were applied before a barrier entered the typology, and both are stated so that a reader
can disagree with them:

\begin{enumerate}
\item \textbf{Dispersion gate.} A barrier is admitted only if it was observed under a task that
      produced non-trivial variance ($SD > 0.42$). Barriers observed only under saturated tasks are
      excluded, because the ease rating cannot discriminate them. This excluded nothing from the theme
      set but changed their relative standing.
\item \textbf{Corroboration annotation.} Every barrier is annotated with its independent code-audit
      status, so that convergence and divergence are visible per barrier rather than as a single
      aggregate claim.
\end{enumerate}

Counts are reported with \textbf{explicit denominators} rather than pooled, because the narrative
corpus and the SEQ corpus cannot be confirmed to be the same ten people (Section~\ref{sec:threats}, T-1).

\section{Findings}

\subsection{The friction typology}

Six themes emerged from analysis. We organise them into five classes, each named for the organisational
capacity it constrains rather than for the interface element where it appears.

\begin{table}[htbp]
\caption{Barrier themes across the two corpora. Navigation and state persistence (8 of 10) and address
capture (7 of 10) dominate; currency appears only in the second SEQ round.}
{\footnotesize
\begin{tabular}{@{}p{0.06\textwidth} p{0.16\textwidth} p{0.15\textwidth} p{0.14\textwidth} p{0.39\textwidth}@{}}
\toprule
\textbf{\#} & \textbf{Theme} & \textbf{Narrative (of 10)} & \textbf{SEQ (of 10)} & \textbf{Observed barrier} \\
\midrule
\textbf{T1} & Currency \& technical register & 3 named (synthesis asserts unanimity) & 2 --- both round 2, \textbf{none in round 1} & USD on an Indonesian site; jargon (``precision diagnostic'', ``internal evaporator'', ``indoor--outdoor unit relation'') over-excludes lay users; technician-internal documentation shown to customers \\
\textbf{T2} & Price transparency & 3 & 1 (+1 currency remark) & Price breakdown not collapsible; unit scope (which AC, indoor or outdoor) unstated \\
\textbf{T3} & Scheduling rigidity & 6 & 5 & Date selected twice with \textbf{disagreeing} options; 3--5 day window read as weekday-only; no calendar; no arrival range \\
\textbf{T4} & Address capture \& validation & 7 & 5 & No map, geocoding, or autocomplete; no saved address; postal code validated but street address not; mandatory fields unmarked; ``service address'' misread as the provider's office \\
\textbf{T5} & Navigation \& state persistence & 6 & 8 & Form data lost on Back/header navigation and on reload; footer links return to the home page; search query persists after clearing; confirmation button stuck in ``processing'' \\
\textbf{T6a/b/c} & Support, confirmation clutter, accessibility & 4 / 1 / 2 & 4 / 1 / 1 & Help buried in footer; three-item FAQ; duplicate WhatsApp and print buttons, two tracking bars, unexplained QR code; type size and grey-on-white contrast below threshold \\
\bottomrule
\end{tabular}}
\end{table}

\begin{table}[htbp]
\caption{Mapping of themes to friction classes.}
\label{tab:classmap}
\begin{tabular}{@{}p{0.24\textwidth} p{0.13\textwidth} p{0.53\textwidth}@{}}
\toprule
\textbf{Class} & \textbf{Themes} & \textbf{Construct} \\
\midrule
\textbf{F1 Service-concept ambiguity} & T2, part of T1 & Information retrieval; faceted-search taxonomy quality \\
\textbf{F2 Infrastructure asymmetry} & T4, part of T3 & Platform affordances; capability vs.\ delivered value \\
\textbf{F3 Representational inconsistency} & T1, T3, T6b & Organisational computing; information consistency \citep{braniwati2026price} \\
\textbf{F4 Transaction-state fragility} & T5, T6a & Process-support systems; transaction persistence \\
\textbf{F5 Support-channel occlusion} & T6c, T8 & Affordance labelling; discoverability \\
\bottomrule
\end{tabular}
\end{table}

\subsection{Five propositions}

We state the typology as propositions rather than as a descriptive list, because the claims are
conditional and the boundary conditions are the substance.

\begin{proposition}[Service-concept ambiguity (F1)]
\label{prop:F1}
\emph{Where a platform's service taxonomy contains near-synonymous entries that are not explicitly
differentiated, users will fail to discriminate among them at a materially higher rate than on an
otherwise identical task differing only in a well-differentiated category --- and the resulting failure
will be misdiagnosed as a navigation or layout problem unless the instrument contains a matched task
pair.}
\end{proposition}

\noindent\textbf{Evidence:} the T1/T2 matched pair (Section~\ref{sec:matched}). \textbf{Boundary
conditions:} the divergence is observable only if tasks are matched on all other properties; an
unmatched instrument cannot distinguish F1 from a general discovery problem.

\begin{proposition}[Infrastructure asymmetry (F2)]
\label{prop:F2}
\emph{Where a transactional platform omits reusable infrastructure --- address normalisation, stored
customer records, order history --- the burden of that omission recurs at every transaction and is borne
by the customer rather than the firm, inverting the economics of repeat custom on which small service
firms depend.}
\end{proposition}

\noindent\textbf{Evidence:} T4 (7 of 10 narrative, 5 of 10 SEQ) with geocoding, autocomplete and saved
addresses as the most-requested features in the study. \textbf{Boundary conditions:} assumes repeat
transaction with a stable customer base, which holds for maintenance services and not for one-off
purchases.

\begin{proposition}[Representational inconsistency (F3)]
\label{prop:F3}
\emph{Where a platform presents two or more representations of the same organisational fact and those
representations disagree, users cannot determine which is authoritative, and in a pricing context they
will treat neither as credible --- regardless of the absolute correctness of either.}
\end{proposition}

\noindent\textbf{Evidence:} the scheduling case (Section~\ref{sec:narrowing}) and the tax-label case
\citep{braniwati2026price}. \textbf{Boundary conditions:} redundancy is benign when copies agree; the
proposition is about disagreement, not duplication.

\begin{proposition}[Transaction-state fragility (F4)]
\label{prop:F4}
\emph{Where a multi-stage transaction does not preserve state across non-linear navigation, a
proportion of completed forms are discarded at a point the user experiences as the end of their effort,
and in a high-commitment service transaction this converts to abandonment or unsupported bookings rather
than retry.}
\end{proposition}

\noindent\textbf{Evidence:} T5 (8 of 10 SEQ, the highest single count in the study) for state loss on
Back and reload. \textbf{Boundary conditions:} we observed state loss, not abandonment. The conversion
from loss to abandonment is \textbf{inferred, not measured} --- no abandonment instrument was in place.
We flag this as the weakest link in our chain and do not build priority 1 on it.

\begin{proposition}[Support-channel occlusion (F5)]
\label{prop:F5}
\emph{Where a service platform provides a support function through a control whose visual form and label
denote a different action, the support channel will be systematically under-used by first-time users,
and the resulting difficulty will be misread as a comprehension deficit rather than a labelling defect.}
\end{proposition}

\noindent\textbf{Evidence:} T8, the lowest-scoring task in the instrument ($M = 5.70$) and the widest
dispersion ($SD = 2.11$, full 1--7 range). Participants found the WhatsApp control and read it as a
\emph{booking} action; some located help only in the footer; one found no direct contact number.
\textbf{Boundary conditions:} the mechanism is labelling and placement, not comprehension --- a
participant who found the channel rated it 7, so training does not substitute for relabelling.

\subsection{F1 in detail: the matched pair isolates a taxonomy defect}
\label{sec:matched}

T1 and T2 are the same task with a different category. T1 (find AC cleaning) reached \textbf{$M = 6.80$,
$SD = 0.42$}. T2 (find leak detection) reached \textbf{$M = 6.20$, $SD = 0.92$} --- a 0.60-point drop
with triple the variance, on an otherwise identical task.

Participants attributed the difficulty to near-synonymous service labels (``smart leak detection''
versus pipe inspection) that the interface did not differentiate, and in two cases to the booking
control sitting below the fold. This is a \textbf{labelling} problem, not a \textbf{layout} problem: the
category was reachable, but participants could not determine \emph{which} entry addressed their problem.

The consequence is disproportionate in cost. Merging two near-synonymous catalogue entries is a
\textbf{data change}. Redesigning the discovery screen is an \textbf{interface change}, an order of
magnitude more expensive. An instrument without a matched pair would report this as ``search is
difficult'' and prescribe the expensive remedy.

\subsection{F3 and the division of labour with code inspection}
\label{sec:division}

The prototype received an independent code-based heuristic audit producing 0 Critical, 2 Major and 4
Minor findings.

\begin{table}[htbp]
\caption{Division of labour between code inspection and participant observation.}
\label{tab:division}
\begin{tabular}{@{}p{0.30\textwidth} p{0.34\textwidth} p{0.26\textwidth}@{}}
\toprule
\textbf{Code-audit finding} & \textbf{Corresponding user-observed barrier} & \textbf{Verdict} \\
\midrule
Major: home-page price differs from catalogue/detail & T2 --- unexplained total jumps; catalogue pricing confusing & \textbf{Converges} \\
Major: no visible keyboard focus indicator & --- & \textbf{Not testable} by this method \\
Minor: header section links fragile & T5 --- footer links returned users to home & \textbf{Converges} \\
Minor: stepper icons contradict text & T6b --- duplicate tracking status bars & \textbf{Converges} \\
Minor: no re-order path & T5 --- participants requested order history & \textbf{Converges} \\
Minor: postal code silently truncated & T4 --- postal code validated, street address not & \textbf{Partly converges} \\
--- & \textbf{T1 --- USD currency on a localised site} & \textbf{Audit missed; testing found} \\
--- & \textbf{Price presentation not independently verifiable} & \textbf{Audit missed; testing found} \citep{braniwati2026price} \\
\bottomrule
\end{tabular}
\end{table}

\textbf{Four of six audited defects were independently rediscovered by participants without
prompting}, which supports heuristic inspection's continued use as a cheap pre-test filter for a firm
with limited evaluation budget. The two misses share a structure: neither is a fault \emph{in} a
control. \texttt{\$} is correctly formatted; a number is correctly formatted. Both are properties of
the \emph{relation between} what the code renders and what a user would expect to follow, and neither
is expressible as a property of any single component.

Because the audit records these defects as \emph{repaired} while the tested sessions reported all of
them, and no session record carries a build identifier, the tested build probably predates the recorded
audit remediation (Section~\ref{sec:threats}, T-2).

\subsection{A narrowing result: disagreement, not duplication}
\label{sec:narrowing}

One finding \textbf{narrows} rather than repeats an earlier complaint, and we give it separate
treatment because a reviewer who misses it would prescribe the wrong fix.

On scheduling (T3), a respondent stated that seeing the date \emph{twice} is acceptable \textbf{if it
functions as a cross-check --- provided the two instances agree}. The defect is therefore
\textbf{disagreement between two date pickers}, not duplication. A single source of truth is a smaller
change than removing one picker.

This is P3 instantiated, and it is the clearest evidence in the study that users distinguish redundancy
from inconsistency. It also illustrates why the remedy-cost classification matters: the naive reading
(``remove the duplicate picker'') prescribes a UI change; the correct reading (``make one authoritative
and derive the other'') prescribes a smaller and more durable fix.

\subsection{What worked}

The prototype's \emph{structure} is sound. Four tasks sat at or near ceiling with almost no dispersion.
The booking flow follows an expected mental model --- browse, detail, schedule, address, pay, confirm
--- and participants navigated it without instruction.

\textbf{Note the asymmetry in T5: the address form reached ceiling ($M = 6.90$, $SD = 0.32$) while
generating the most frequently requested features in the entire study.} The form is easy to fill and
wrong in design. Ease of completion and fitness for purpose are independent properties, and SEQ
measures only the first. For a service firm, this asymmetry is dangerous in a specific way: the
address form looks like the platform's strongest flow in any dashboard built on ease metrics, while
being the point at which the customer does the most avoidable work.

\subsection{A contradiction between our own sources, and how we resolved it}

The narrative source report's recommendation matrix rates support/contact remediation \textbf{Low}
priority. Both recap sheets rate it \textbf{High}. The narrative's own session data contradicts its own
prioritisation, since four of its sessions record difficulty locating support.

We adopt the recap sheets' rating, on two grounds: it is supported by the largest weight of recorded
evidence, and it is supported by the dispersion analysis, which places T8 lowest in mean ease and
highest in variance in the entire instrument --- the profile of a barrier that is real rather than an
artefact of a saturated task.

We report the disagreement rather than suppressing it, for two reasons. A reader checking our
recommendation against the primary documents should be able to see that we noticed. And the
disagreement is itself \textbf{evidence about a failure mode of multi-source research projects} (R4):
where a narrative synthesis and quantitative recap sheets are produced separately, their prioritisation
can diverge, and the synthesis is the artefact most likely to be read as authoritative. \textbf{We
recommend that where source records disagree on priority, the disagreement is recorded and resolved
explicitly against recorded evidence rather than silently in favour of the narrative conclusion.}

\section{Remediation}
\label{sec:remed}

\subsection{Remedy-cost classification}

The central practical contribution. Each finding is classified by the layer at which it must be
remedied, because the layers differ by an order of magnitude in cost and in durability.

\begin{table}[htbp]
\caption{Remedy-cost classification by layer.}
\label{tab:layers}
\begin{tabular}{@{}p{0.16\textwidth} p{0.26\textwidth} p{0.14\textwidth} p{0.36\textwidth}@{}}
\toprule
\textbf{Layer} & \textbf{Meaning} & \textbf{Indicative relative cost} & \textbf{Findings} \\
\midrule
\textbf{D --- Data change} & Catalogue content, labels, config.\ No code. & Lowest & F1 (merge near-synonymous entries); F3 (single date source of truth); F5 (relabel support control); register (remove technician-internal text) \\
\textbf{I --- Interface change} & Layout, validation, routing, persistence within the app. & Moderate & F4 (persist checkout state); F3 (collapse confirmation clutter); F5 (move support to header); F2 (mark required fields) \\
\textbf{N --- Infrastructure provision} & Requires platform capability the firm cannot add. & Highest & F2 (geocoding, address autocomplete, stored addresses, order history, re-order, calendar) \\
\bottomrule
\end{tabular}
\end{table}

\textbf{This distinction is invisible in standard usability reporting}, which lists defects without
reference to the layer that must change. A triage meeting reading a flat list will plausibly fix the
label on a support button and redesign the address form, when the reverse order is indicated: the data
change is cheaper, faster, and does not degrade with the next release.

\subsection{Evidence-linked remediation matrix}

Every row traces to a frequency count, a dispersion measure, and a code-audit status. Priority is
auditable, not asserted.

{\footnotesize
\begin{longtable}{@{}p{0.03\textwidth} p{0.20\textwidth} p{0.05\textwidth} p{0.24\textwidth} p{0.07\textwidth} p{0.07\textwidth} p{0.05\textwidth} p{0.07\textwidth} p{0.06\textwidth}@{}}
\caption{Evidence-linked remediation matrix. Every row traces to a frequency count, a dispersion measure, and a code-audit status.}\\
\label{tab:matrix}
\toprule
\textbf{P} & \textbf{Finding} & \textbf{Class} & \textbf{Remedy} & \textbf{Layer} & \textbf{Freq.} & \textbf{SD} & \textbf{Audit} & \textbf{Sev} \\
\midrule
\endfirsthead
\multicolumn{9}{@{}l}{\footnotesize\itshape Table \thetable\ continued from previous page}\\
\toprule
\textbf{P} & \textbf{Finding} & \textbf{Class} & \textbf{Remedy} & \textbf{Layer} & \textbf{Freq.} & \textbf{SD} & \textbf{Audit} & \textbf{Sev} \\
\midrule
\endhead
\midrule
\multicolumn{9}{r@{}}{\footnotesize Continued on next page}\\
\endfoot
\bottomrule
\endlastfoot
\textbf{1} & Price presentation not independently verifiable (ambiguous tax base) & F3 & State the tax base in the label; assert \texttt{amount == rate x base} at build time & \textbf{D} + pipeline & 3 & --- & Missed by audit & \textbf{Critical} \\
\textbf{2} & USD on an Indonesian site; technician-internal jargon in customer-facing text & F3 & Supply IDR pricing; rewrite the service taxonomy in operational lay register; remove internal procedure text & \textbf{D} & 3 & 1.27 & Missed by audit & \textbf{High} \\
\textbf{3} & Support channel undiscoverable; control labelled as booking & F5 & Persistent, explicitly labelled support action in the header, not the footer; add direct contact; FAQ search and per-service FAQ & \textbf{D} + \textbf{I} & 4 & \textbf{2.11} & --- & \textbf{High} \\
\textbf{4} & Form state lost on Back/reload & F4 & Persist checkout state; warn before destructive navigation; offer resume of an incomplete booking & \textbf{I} & 6 & 0.95 & Converges & \textbf{High} \\
\textbf{5} & Street address unvalidated; mandatory fields unmarked; ``service address'' misread as the firm's office & F2 & Validate street address on the same terms as postal code; mark required/optional; relabel as home address & \textbf{I} & 7 & 0.32 & Partly converges & \textbf{High} \\
\textbf{6} & Two date pickers that \textbf{disagree} & F3 & Single source of truth for available dates; one picker; retain a read-only cross-check guaranteed identical & \textbf{D} & 6 & 0.42 & --- & \textbf{High} \\
\textbf{7} & Near-synonymous catalogue entries (T1/T2 matched pair) & F1 & Merge the entries; differentiate in the label, not only in the description & \textbf{D} & --- & 0.92 & --- & \textbf{High} \\
\textbf{8} & Footer links return to home; search query persists after clear & F4 & Fix footer routing; reset query on clear and navigation & \textbf{I} & 6 & 0.95 & Converges & \textbf{Medium} \\
\textbf{9} & Confirmation-page clutter: duplicate buttons, two tracking bars, unexplained QR & F3 & Collapse to a static transaction summary; move tracking to order history; label or remove the QR & \textbf{I} & 1 & 1.27 & Converges & \textbf{Medium} \\
\textbf{10} & Type size and contrast below threshold & F3 & Minimum 12 pt; secondary text to WCAG AA; verify rendered contrast in a browser & \textbf{I} & 2 & 1.27 & Flagged, not testable & \textbf{Medium} \\
\textbf{11} & No geocoding, saved address, order history, re-order & F2 & Map pin selection; address autocomplete; addresses stored by label; order history with re-order & \textbf{N} & 7 & 0.32 & Converges & \textbf{Medium} \\
\textbf{12} & No path for a missing service; rigid scheduling window & F2 & ``Other service'' request path; full calendar; arrival-time range; custom-time request & \textbf{N} + \textbf{I} & 6 & 0.42 & --- & \textbf{Low} \\
\end{longtable}}

\textbf{On priority 1.} We downgrade the price-integrity item from the \emph{Critical: total $\neq$ sum
of its own line items} of earlier drafts. A reproduction attempt against the prototype source
\textbf{refuted} that finding: no build produces an inconsistent total, and the arithmetic the observed
figures were checked against was itself computed on a different tax base than the code implements
\citep{braniwati2026price}. The item survives at Critical only in its verified form --- an ambiguous
tax base (Section~\ref{sec:division}), a labelling defect that makes the displayed amount
unverifiable against the displayed rate under its most natural reading. \textbf{We report this change
explicitly because a reader comparing against the earlier version is entitled to know a severity was
revised and why.}

\subsection{Implications}

\textbf{For MSME service firms.} Evaluate the interface as an organisational instrument before launch.
The diagnostic that mattered most in this study --- two pickers that disagree, a control that means
something other than it says --- is invisible to the owner of the business and to any ease metric. A
one-hour review in which the owner attempts each task against a written criterion, and states aloud
what each screen claims, would have caught priority 1, 3 and 6.

\textbf{For platform designers serving MSME sectors.} Ship the D-layer fixes first. In this dataset the
four highest-frequency barriers had D-layer components, and three of them had \emph{only} D-layer
components. The finding is uncomfortable but actionable: a disproportionate share of
organisational-credibility risk in MSME marketplaces lives in catalogue content, not in interface code.

\textbf{For researchers.} Build matched task pairs into formative instruments where taxonomies are
involved. The T1/T2 contrast cost one additional task and converted an unresolvable complaint into a
priced remedy. Where designers cannot add probes to a saturated instrument, at minimum add one defect
probe: on a competent design it is the only task type that does not saturate
\citep{braniwati2026dispersion}.

\textbf{On the evidence base.} All findings derive from one prototype, one market and ten purposively
recruited participants. The typology is offered as a starting vocabulary for comparable settings, with
its propositions stated conditionally so that a later study can test or refute them. \textbf{We do not
claim prevalence in any market.}

\section{Ethical considerations}
\label{sec:ethics}

Written informed consent was obtained from all ten participants before participation began. The
instrument is reproduced in Appendix~\ref{app:consent}. Participants were informed of the study's
purpose, that the prototype simulated every transaction with no payment, booking, technician dispatch
or contact, that a single fictional address was used for all participants, and that they could decline
any question or stop at any point without consequence. No names, telephone numbers, home addresses,
passwords or payment data were collected or retained; all session records are held under pseudonymous
codes (P01--P10), and re-identification from the analysis files is not possible.

Screen and audio recording was a separate opt-in that all ten participants accepted. The recordings are
held under institutional control --- in BINUS OneDrive and a local copy, accessible only to the research
team --- and are scheduled for deletion on \textbf{1 January 2027}. They were not used in this analysis.

Two features of that arrangement are disclosed rather than omitted. \textbf{First}, the
recording-policy fields in the approved consent instrument --- what is recorded, purpose, who may
access, storage location, deletion date, and the contact for deletion requests --- were completed
before recording began and given to every participant who consented. \textbf{Second}, no additional
encryption and no backup were applied; loss of the files before the deletion date would therefore be
permanent. A determination covering the analysis and publication stage has been requested from the
BINUS Research Ethics Committee.

\textbf{No institutional research ethics review was conducted prior to data collection.} This is a
departure from our own study protocol, which specified ethics determination as a precondition for
commencing formal sessions, and from BINUS University's requirement that research involving human
participants be submitted for determination before recruitment. We disclose it as a limitation. The
safeguards described above were applied in every session and are evidenced by the signed consent
instruments held by the corresponding author. \textbf{We note that these protections are not
equivalent to independent ethical review and do not substitute for it}, and we recommend that
subsequent work in this area obtain prospective determination. The consent instrument was an
internally drafted form that was not ratified by an institutional ethics committee; its text is
provided in full in Appendix~\ref{app:consent} so that readers can assess it directly.

A determination has been requested from the BINUS Research Ethics Committee under the exemption track
for secondary analysis of deidentified data. \textbf{[If granted: This study was determined to require
no ethical review by the Research Ethics Committee, BINUS University (determination no.~[X],
[date]).]} \textbf{[If refused or pending: No determination has been issued as of submission.]} The
request, and a full disclosure of the departure, are available in the study's repository.

\section{Threats to validity}
\label{sec:threats}

\textbf{T-1 --- Participant rosters do not agree across source records. [Severe]} The narrative corpus
names ten participants; the two quantitative recap sheets name a partly different set --- four names
overlap, three appear only in the recap sheets, and six narrative participants have no SEQ record
anywhere. Both recap sheets end with a note that missing names be filled in. \textbf{Consequence:} the
narrative and SEQ corpora cannot be confirmed to be the same ten people. We never aggregate across them
as though they were, and every count in Section~\ref{tab:classmap} carries an explicit denominator.

\textbf{T-2 --- No build identifier, date, or moderator record. [Severe]} The protocol required recording
the canonical prototype and its build hash before sessions began; no session record contains one. As
Section~\ref{sec:division} shows, the tested build probably predates the recorded audit remediation.
\textbf{Consequence:} no result may be attributed to the current prototype.

\textbf{T-3 --- Demographic detail minimal. [Low]} Participants were Gen Z, aged 19--21, verified against
national ID cards at recruitment; no date of birth was transcribed into the research data, and gender
and occupation were not collected. \textbf{Consequence:} the sample's composition is established only
to the extent stated, and the organisational claims in Section~\ref{sec:remed} are correspondingly
bounded to what a purposive sample of ten young adults can support.

\textbf{T-4 --- Ethics and consent. [Severe]} Written informed consent was obtained from all ten
participants and the safeguards applied in every session are documented
(Section~\ref{sec:ethics}). \textbf{No institutional research ethics review was conducted prior to data
collection}, a departure from our own protocol and from BINUS's requirement that research involving
human participants be submitted for determination before recruitment. A determination has been
requested from the Research Ethics Committee under the exemption track for secondary analysis of
deidentified data. \textbf{Until a written determination is issued, this threat remains open}, and
readers should treat the dataset's provenance as incomplete. We note that consent and institutional
review are separate requirements, and that satisfying the first does not discharge the second.

\textbf{T-5 --- The abandonment inference is not measured. [Moderate]} P4 asserts that state loss
converts to abandonment or unsupported bookings. We observed state loss; \textbf{we did not observe
abandonment}, and no instrument would have detected it. The conversion is an inference from the
structure of the transaction, and a reader who rejects it should reject priority 4's severity.

\textbf{T-6 --- Ceiling effect limits four tasks. [Moderate]} T1, T3, T4, T5 ($SD \leq 0.42$) cannot
discriminate between design states. Later iterations will need harder tasks, a more complex stimulus,
or different measures.

\textbf{T-7 --- Variance collapse between rounds. [Moderate]} Round-2 per-task $SD$ (0.49) is
$2.8\times$ smaller than round-1 (1.39), with three participants returning near-uniform maximum
scores. Possible causes --- differing scoring discipline, moderator effects, demand characteristics,
transcription convention --- cannot be distinguished from the available records. We draw \textbf{no}
design-effect conclusion; the build did not change between rounds.

\textbf{T-8 --- Qualitative corpus is observer paraphrase, not transcript. [Moderate]}
\textbf{Consequence:} no direct participant quotation is used anywhere in this paper.

\textbf{T-9 --- Single study, single domain. [Moderate]} One prototype, one market, one service
category. The propositions are stated conditionally and are offered as testable claims, not as
established general theory.

\textbf{T-10 --- References partly unverified. [Low]} See the study repository provenance record.

\section{Conclusion}

We evaluated an Indonesian-language home-services marketplace prototype built for micro and small
service firms, with ten participants across eight tasks and two rounds, cross-validated against an
independent code audit. Aggregate task success was 93.75\%. The aggregate is close to uninformative,
and the informative structure beneath it is not a list of interface defects but a \textbf{typology of
five friction classes}: service-concept ambiguity, infrastructure asymmetry, representational
inconsistency, transaction-state fragility, and support-channel occlusion.

The organisational reading changes what the list means. These are not usability problems to be
incrementally smoothed. They are constraints on a micro-firm's capacity to be trusted by a stranger
before any interaction has occurred, and in a transaction where the customer cannot substitute a
provider and cannot independently check the quotation, the cost of that constraint falls on the
customer.

Three practical results follow. \textbf{Remediation cost is not proportional to user effort}, and three
of the four highest-frequency barriers had remedies available entirely in the catalogue-data layer ---
no code, no redesign. \textbf{A matched task pair distinguishes a taxonomy defect from a layout
defect}, converting an unresolvable complaint into a data change. And \textbf{ease and fitness are
independent properties}: the address form topped the ease scale while generating the study's
most-requested features, which means any remediation dashboard built on ease metrics would have shown
the platform's weakest point as its strongest.

We offer the typology as testable propositions rather than established theory, and we report the five
provenance defects that constrain them --- including one source-record contradiction we resolved
against its own narrative, and one severity we revised downward after a reproduction refuted it. A
finding that survives its own refutation attempt is worth more to a practice than one that was never
tested.

\section*{Acknowledgement(s)}

Meutia Braniwati chaired the study and contributed the visual branding analysis; Aditya Paramananda and
Tan, Michael Wilson Setiawan designed the prototype; Stephanie Evangelina Pramudito handled
administration and visual design. Their contributions are acknowledged here; they are not authors of
this paper.

\section*{Disclosure statement}

No potential conflict of interest was reported by the authors.

\section*{Funding}

No external funding was received.

\section*{Notes on contributor(s)}

Stated in CRediT terms. Bagaskoro Saputro: Methodology, Investigation, Formal analysis, Writing ---
original draft. Yashella Tirana: Formal analysis (usability-test analysis, heuristic evaluation),
Writing --- review and editing. Freysia Chandra Saliman: Investigation, Data curation, Resources,
Writing --- review and editing. Laurentius Ryan Dofani Pradita, Gregory Adrianus Sugiono, Gregorius
Andrew Rahardja, Steven Karunia Liem, and Joshua Joyson Kustiadi: Investigation (data collection;
usability-test moderation), Data curation (data annotation). All authors read and approved the final
manuscript. Bagaskoro Saputro is responsible for the study design and for the findings reported here.

\section{Data availability statement}

All data underlying every finding in this paper will be deposited on OSF at
\textbf{[OSF\_LINK --- to be added]}. The deposit contains the friction-coding sheets behind the
typology (Section~\ref{tab:classmap}), the heuristic evaluation records, the evidence-linking table for
the remediation matrix (Table~\ref{tab:matrix}), the theme-to-class comparison (Appendix~\ref{app:themeclass}),
and the provenance and data-management record (Appendix~\ref{app:prov}), all under the pseudonymous
participant codes P01--P10 used here.

The participant information and consent instrument is reproduced in Appendix~\ref{app:consent}. Screen
and audio recordings are \textbf{not} shared: they were not used in this analysis and are scheduled for
deletion on 1 January 2027 under the approved consent instrument.

\section{References}

\bibliographystyle{apalike}
\bibliography{paper-3-jocec-references}

\appendix

\section{Comparison of the six themes against the five friction classes}
\label{app:themeclass}

\begin{table}[htbp]
\caption{Theme-to-class comparison.}
{\footnotesize
\begin{tabular}{@{}p{0.22\textwidth} p{0.08\textwidth} p{0.22\textwidth} p{0.20\textwidth} p{0.18\textwidth}@{}}
\toprule
\textbf{Theme} & \textbf{Class} & \textbf{Construct} & \textbf{Primary evidence} & \textbf{Independent code finding} \\
\midrule
T1 Currency \& register & F3 & Information consistency & 3 narr / 2 SEQ (both round 2) & Missed \\
T2 Price transparency & F1 & IR taxonomy quality & 3 narr / 1 SEQ & Converges (home-page price) \\
T3 Scheduling rigidity & F2 + F3 & Capability + consistency & 6 narr / 5 SEQ & --- \\
T4 Address capture & F2 & Platform affordances & 7 narr / 5 SEQ & Partly converges (postal code) \\
T5 Navigation \& state & F4 & Process-support systems & 6 narr / 8 SEQ & Converges (links, re-order) \\
T6a Confirmation clutter & F4 & Process-support systems & 1 narr / 1 SEQ & Converges (stepper icons) \\
T6b Duplicate tracking & F3 & Information consistency & 1 narr / 1 SEQ & Converges (stepper icons) \\
T6c Accessibility & F3 & Interface legibility & 2 narr / 1 SEQ & Flagged, not testable \\
T8 Support findability & F5 & Affordance labelling & 4 narr / 4 SEQ, SD 2.11 & --- \\
\bottomrule
\end{tabular}}
\end{table}

\noindent T6a/b/c were counted as three sub-themes in the source report; the class mapping is ours, and
T8 was analysed as a task-level finding rather than a theme in the source report. We present it as a
theme because it satisfies the dispersion gate and carries the study's highest-variance task.

\section{Provenance and data management}
\label{app:prov}

Full traceability record --- source inventory, roster cross-comparison, recomputed statistics, and the
list of figures permanently barred from this study's reporting --- is maintained in the study
repository.

\textbf{Prohibited figures.} The superseded draft contained 87.5\%, 4.20/5, 84.6\%, and 4.23/5. These
were never derivable from any source and remain permanently barred. Note that 84.6\% is \textbf{not}
the corrected completion rate (93.75\%) and 87.5\% is \textbf{not} any figure in
Section~\ref{sec:method}; the resemblance is coincidental and must not be used to rehabilitate them.

\textbf{Severity revision record.} Priority 1 was \emph{Total inconsistent with its own line items}
(Critical) in earlier drafts. A reproduction attempt against the prototype source refuted that finding.
The item is restated as an ambiguous tax base, a labelling defect verified from source, and retained at
Critical only in that verified form. See \citet{braniwati2026price}.

\section{Participant information and consent instrument}
\label{app:consent}

Reproduced in full so that readers can assess the instrument directly rather than take the authors'
characterisation of it on trust. This is the form given to participants; it carries the internal status
marking \emph{Draf.\ Belum disetujui}, disclosed in Section~\ref{sec:ethics}. The Part 3
recording-policy fields were \textbf{completed and approved before use}, and their approved values are
reproduced below. The blank master is maintained in the study repository; the signed copies, bearing
participant signatures, are held by the corresponding author in institutional storage and are
deliberately \textbf{not} included in the repository.

\subsection*{Part 1 --- Information for participants}

\textbf{What is this research?} We are testing a website for finding and booking home-maintenance
services, for example air-conditioner cleaning or leak repair. We want to know whether the website is
easy to understand and whether its information is clear enough.

\textbf{What are we testing?} We are testing the website's design, not your ability. There are no right
or wrong answers.

\textbf{What will you do?} Read and agree to this information; complete some simple tasks on the
website while saying what you are thinking; answer brief questions afterwards.

\textbf{What is not part of this research?} There is \textbf{no payment} --- you will not be billed and
will not be asked for card details. There is \textbf{no real booking} --- the confirm button completes
a display simulation only. \textbf{No technician is dispatched} and \textbf{no telephone number is
dialled}. There is \textbf{no GPS tracking} --- the tracking page is an illustration, not a real
location. \textbf{No account is created}, and any password you type is not stored.

\textbf{About prices and other information.} Some information in the prototype is deliberately
unfinished. Displayed prices are sample figures used to test the display. Names, reviews and technician
profiles are also sample data. Service booking is not connected to any system.

\textbf{Will I be asked for personal data?} No. We do not ask for your full name, telephone number,
home address, password, or payment data. For the booking task we give you a fictional address, the
same one for every participant.

\textbf{Is the session recorded?} Recording takes place \textbf{only if you agree to it}. Recording is a
separate option and you may refuse it without affecting your participation. Details of who may access
recordings and when the data is deleted are provided in our official form.

\textbf{What are my rights?} You may skip any question. You may stop at any time without giving a
reason. Stopping will not affect your grade, the study, or your relationship with the university or its
partners. You may request that your session data be withdrawn under the applicable procedure.

\textbf{What happens to the results?} Notes are analysed by the research team and reported anonymously.
Any quotations do not include your name. Results are used to improve the website's design.

\textbf{Confirmation of understanding.} $\square$ I have read and understood this information.
$\square$ I have had the opportunity to ask questions and receive answers. $\square$ I understand what I
am being asked to do. $\square$ I understand that there is no payment, booking, or real contact.

\subsection*{Part 2 --- Consent to participate}

Complete only after the spoken explanation has been given. Use the participant code, not the name, on
analysis documents.

\begin{center}
\begin{tabular}{@{}p{0.78\textwidth}cc@{}}
\toprule
\textbf{Statement} & \textbf{Yes} & \textbf{No} \\
\midrule
I am at least 18 years old. & $\square$ & $\square$ \\
I understand the purpose, activities and duration of the session. & $\square$ & $\square$ \\
I understand this session involves no payment, booking, or real contact. & $\square$ & $\square$ \\
I understand I may skip questions or stop at any time. & $\square$ & $\square$ \\
I understand session data will be stored under a code and reported anonymously. & $\square$ & $\square$ \\
\textbf{I agree to take part in this session.} & $\square$ & $\square$ \\
\bottomrule
\end{tabular}
\end{center}

\textbf{Participant's decision:} $\square$ agree to take part \quad $\square$ do not agree to take part

Participants may refuse without giving a reason. The moderator records refusals as a count, not as
content, and it does not affect subsequent recruitment.

\subsection*{Part 3 --- Consent to recording (separate)}

Complete only if recording is proposed in the protocol. Recording is not required, and refusing to be
recorded does not invalidate consent to participate.

\begin{center}
\begin{tabular}{@{}p{0.78\textwidth}cc@{}}
\toprule
\textbf{Statement} & \textbf{Yes} & \textbf{No} \\
\midrule
I agree to be recorded, screen and/or audio, during the session. & $\square$ & $\square$ \\
I agree my recording may be used for internal research purposes. & $\square$ & $\square$ \\
I agree the recording may be reused for publication in anonymised form. & $\square$ & $\square$ \\
\bottomrule
\end{tabular}
\end{center}

\begin{center}
\begin{tabular}{@{}p{0.30\textwidth}p{0.50\textwidth}@{}}
\toprule
\textbf{Policy item} & \textbf{Approved value} \\
\midrule
What is recorded (screen, audio, or both) & Screen and audio \\
Purpose of recording & Observer-note transcription and verification of the analysis \\
Who may access it & The research team \\
Storage location & BINUS OneDrive (institutional account) and a local copy \\
Deletion date & 1 January 2027 \\
Contact for deletion requests & Corresponding author \\
\bottomrule
\end{tabular}
\end{center}

\noindent These values were approved before recording began and were given to every participant who
consented to recording. All ten participants consented.

\noindent Recording must not be performed on a build that requests payment, collects real addresses, or
accesses GPS.

\subsection*{Part 4 --- Consent to quotation (separate)}

\begin{center}
\begin{tabular}{@{}p{0.78\textwidth}cc@{}}
\toprule
\textbf{Statement} & \textbf{Yes} & \textbf{No} \\
\midrule
I agree my quotations may be used anonymously in the report and in publications. & $\square$ & $\square$ \\
\bottomrule
\end{tabular}
\end{center}

\noindent Where a participant permits quotation, still use the code \texttt{P\_\_}. Do not link a
participant code to a name in the same document. Quotations may be used only where they come from notes
or recordings that actually exist. \textbf{Do not rewrite a moderator's note into a direct quotation.}

\subsection*{Part 5 --- Pre-session checklist}

The moderator confirms every item below before the participant begins.

\noindent $\square$ Spoken explanation given and understanding confirmed \quad $\square$ Consent form
signed in Part 2 \quad $\square$ Part 3 recording consent explicitly completed or declined \quad
$\square$ Part 4 quotation consent explicitly completed or declined \quad $\square$ Fictional practice
address prepared and no participant address collected \quad $\square$ Prototype build and timestamp
recorded \quad $\square$ Build verified with \texttt{npm run lint} and \texttt{npm run build} \quad
$\square$ Moderator states that tasks stop if the prototype requests sensitive data or opens payment
\quad $\square$ The way to end the session has been communicated.

\end{document}
